% =====================================================================
%  1bat-ud5-7.tex  —  SESSIÓ 7: La fórmula de l'interès compost
%  (sense preàmbul: s'inclou des de main.tex)
% =====================================================================
\horatitol{Sessió 7 — La fórmula de l'interès compost}%
{De l'índex de variació a la potència: si cada període el capital es multiplica per (1+r), aplicar-lo t períodes és elevar el factor a t.}

\subsection{Idea clau}

A la sessió 6 vam distingir interès simple i compost. Avui formalitzem el
\textbf{compost}, el que de debò mou els estalvis i els deutes, connectant-lo amb una
idea que ja coneixem: l'\textbf{índex de variació}. Si cada any el capital es
multiplica per $(1+r)$, aplicar-ho diversos anys vol dir \textbf{elevar} aquest factor
a una potència.

\begin{idea}
\textbf{La fórmula de l'interès compost}
\[
\boxed{\,C_f = C_0\per (1+r)^{t}\,}
\]
\begin{itemize}[leftmargin=1.4em, itemsep=2pt]
  \item $C_0$ — \textbf{capital inicial} (els diners de partida).
  \item $r$ — \textbf{tipus d'interès per període}, com a factor ($5\%\rightarrow r=0,05$).
  \item $t$ — \textbf{nombre de períodes} (l'exponent).
  \item $C_f$ — \textbf{capital final}.
\end{itemize}
\textbf{La idea clau:} cada període multipliquem per $(1+r)$; fer-ho $t$ vegades és
$(1+r)^t$. Per això l'efecte \textbf{es dispara} com més gran és $t$: és la potència
treballant a favor (estalvi) o en contra (deute).
\end{idea}

\subsection{Exemples}

\begin{exemple}[l'estalvi que creix]
Ingressem $\num{1000}\eur$ al $3\%$ compost anual. Cada any es multiplica per $1,03$:
\[
\begin{aligned}
\text{any 1: } &\num{1000}\per 1,03=\num{1030}\eur,\\
\text{any 2: } &\num{1030}\per 1,03=\num{1060,9}\eur \;=\;\num{1000}\per 1,03^2,\\
\text{any 3: } &\num{1000}\per 1,03^3\approx\num{1092,73}\eur.
\end{aligned}
\]
Fixa't: al $2$n any l'interès ja no és $\num{30}\eur$ sinó $\num{30,9}\eur$, perquè
els $\num{30}\eur$ del primer any \textbf{també generen interès}. Amb interès simple
tindríem només $\num{1090}\eur$ als $3$ anys; amb compost, $\num{1092,73}\eur$.
\end{exemple}

\begin{exemple}[el deute que s'infla]
El mateix mecanisme, en contra. Un deute de $\num{300}\eur$ a un crèdit ràpid del
$10\%$ \textbf{mensual} compost, si no es paga:
\[
\begin{aligned}
\text{mes 3: } &\num{300}\per 1,10^3\approx\num{399,3}\eur,\\
\text{mes 6: } &\num{300}\per 1,10^6\approx\num{531,47}\eur,\\
\text{mes 12: } &\num{300}\per 1,10^{12}\approx\num{941,53}\eur.
\end{aligned}
\]
En un any, el deute s'ha \textbf{triplicat} llarg (de $\num{300}$ a gairebé
$\num{942}\eur$). Aquesta és \textbf{l'espiral} del crèdit abusiu: la mateixa potència
que fa créixer un estalvi modest fa créixer un deute fins a fer-lo impagable.
\end{exemple}

\subsection{Exercicis resolts}

\begin{exresolt}[capital final amb interès compost]
Quant hi haurà al cap de $4$ anys si ingressem $\num{2000}\eur$ al $5\%$ compost anual?
\solucio
Apliquem la fórmula amb $C_0=\num{2000}$, $r=0,05$, $t=4$:
\[
C_f=\num{2000}\per (1,05)^4=\num{2000}\per 1,21550625\approx\num{2431,01}\eur.
\]
\textbf{Interpretació:} els interessos generats han estat uns $\num{431}\eur$, més que
els $\num{400}\eur$ que donaria l'interès simple ($4\per 0,05\per\num{2000}$), gràcies
a l'efecte compost.
\resposta{$C_f\approx\num{2431,01}\eur$.}
\end{exresolt}

\begin{exresolt}[comparar simple i compost]
Una família deu $\num{1000}\eur$ al $4\%$ anual durant $3$ anys. Quant deurà amb
interès \textbf{simple} i quant amb \textbf{compost}?
\solucio
\textbf{Simple:} interès anual $0,04\per\num{1000}=\num{40}\eur$; en $3$ anys
$\num{120}\eur$; deute final $\num{1120}\eur$.

\textbf{Compost:} $C_f=\num{1000}\per (1,04)^3=\num{1000}\per 1,124864\approx\num{1124,86}\eur$.

La diferència ($\num{1124,86}-\num{1120}=\num{4,86}\eur$) encara és petita en $3$ anys,
però creix molt de pressa com més llarg és el termini i més alt el tipus.
\resposta{Simple: $\num{1120}\eur$. Compost: $\approx\num{1124,86}\eur$.}
\end{exresolt}

\subsection{Exercicis per fer}

\noindent\textit{Usa la fórmula $C_f=C_0\per (1+r)^t$. Pots usar la calculadora per a
les potències.}

\begin{exercicis}
\begin{exlist}
  \item Calcula el capital final amb interès compost:
    \begin{enumerate}[label=\alph*), leftmargin=1.6em, topsep=2pt]
      \item $\num{500}\eur$ al $2\%$ anual durant $3$ anys;
      \item $\num{1500}\eur$ al $3\%$ anual durant $5$ anys.
    \end{enumerate}
  \item Un estalvi de $\num{1000}\eur$ al $4\%$ compost anual. Quant hi haurà al cap de
        $1$, $2$ i $10$ anys? Observa com creix la diferència.
  \item Un deute de $\num{500}\eur$ al $8\%$ compost anual que no es paga. Quant es
        deurà al cap de $5$ anys?
  \item \textbf{(Comparació)} Sobre $\num{1000}\eur$ al $5\%$ durant $3$ anys, calcula
        el capital final amb interès \textbf{simple} i amb \textbf{compost}. Quina és
        la diferència?
  \item \textbf{(Interpretació)} En la fórmula $C_f=C_0(1+r)^t$, què passa amb el
        capital final si dupliquem el temps $t$? I si dupliquem el tipus $r$? Per què el
        temps té un efecte tan fort (pensa que és l'exponent)?
  \item \textbf{(Raonament)} Una persona et diu: «el banc m'ofereix un $2\%$ pels meus
        estalvis, però el crèdit ràpid em cobra un $10\%$ mensual; total, són
        percentatges petits». Explica-li, amb el que saps de l'interès compost, per què
        aquesta comparació és molt enganyosa.
\end{exlist}
\end{exercicis}

% ---------------------------------------------------------------------
\fullsolucions{Sessió 7}
\begin{solucions}
\begin{sollist}
  \item (a) $\num{500}\per (1,02)^3\approx\num{500}\per 1,061208=\num{530,60}\eur$.
        \quad (b) $\num{1500}\per (1,03)^5\approx\num{1500}\per 1,159274=\num{1738,91}\eur$.
  \item Any $1$: $\num{1000}\per 1,04=\num{1040}\eur$. Any $2$:
        $\num{1000}\per 1,04^2=\num{1081,6}\eur$. Any $10$:
        $\num{1000}\per 1,04^{10}\approx\num{1480,24}\eur$. La diferència respecte de
        l'interès simple ($\num{1400}\eur$ als $10$ anys) ja és notable: uns
        $\num{80}\eur$ més pel compost.
  \item $\num{500}\per (1,08)^5\approx\num{500}\per 1,469328=\num{734,66}\eur$.
  \item Simple: $\num{1000}+3\per 0,05\per\num{1000}=\num{1150}\eur$. Compost:
        $\num{1000}\per (1,05)^3\approx\num{1157,63}\eur$. Diferència:
        $\num{1157,63}-\num{1150}=\num{7,63}\eur$.
  \item Si dupliquem $t$, el capital final \emph{no} es duplica: es multiplica per
        $(1+r)^t$ \textbf{una altra vegada}, és a dir, l'efecte és molt més que el
        doble (creixement exponencial). Si dupliquem $r$, l'efecte també creix, però el
        temps, com que és l'\textbf{exponent}, fa créixer el capital molt més de pressa
        que un canvi equivalent en el tipus: per això el termini llarg és tan
        determinant.
  \item Resposta oberta. Idea: el $2\%$ del banc és \textbf{anual}, però el $10\%$ del
        crèdit és \textbf{mensual} i, a més, \textbf{compost}: en un any el deute es
        multiplica per $1,10^{12}\approx 3,14$, és a dir, creix més d'un $200\%$, mentre
        que l'estalvi creix només un $2\%$. Comparar «$2\%$» amb «$10\%$» com si fossin
        del mateix tipus amaga que un és anual i petit i l'altre, mensual i devastador.
\end{sollist}
\end{solucions}
